\section*{Chapitre \ref{ch:in:choosing} --- Choix}

\begin{solution}{exo:in:choosing:1}
L'offre de l'ingénieur en technologie, comparée à celle du quant de banque~: une rémunération espérée plus élevée, une fourchette plus étroite, une part moindre prélevée
par l'impôt et le logement, pas d'astreinte, et un apprentissage et une sécurité équivalents.
\end{solution}

\begin{solution}{exo:in:choosing:2}
Rémunération espérée~0.30, risque de rémunération~0.15, lieu~0.10, astreinte~0.15, apprentissage~0.20, sécurité~0.10.
\end{solution}

\begin{solution}{exo:in:choosing:3}
L'offre du teneur de marché sur les deux~: un équivalent certain de 1\,365\,186~\$ et une rémunération espérée de 1\,949\,485~\$ sur
cinq ans.
\end{solution}

\begin{solution}{exo:in:choosing:4}
Elle a la fourchette de rémunération la plus large et la seule chance d'un emploi fermé, les plus mauvais scores sur deux attributs, et elle est devancée
par le teneur de marché sur la rémunération espérée et sur le lieu~; seuls les tirages qui pondèrent fortement l'apprentissage et l'absence d'astreinte, et
faiblement le risque de rémunération et la sécurité, la placent en tête.
\end{solution}

\begin{solution}{exo:in:choosing:5}
Il traite toutes les pondérations comme également probables, y compris des pondérations extrêmes qu'aucun vrai décideur ne retient, et
ignore ce que la personne qui choisit sait de ses propres priorités~; les acceptabilités décrivent la robustesse des offres, et non sa préférence.
\end{solution}

\begin{solution}{exo:in:choosing:6}
Toujours l'offre du teneur de marché (0.894, contre 0.706 pour celle de la plateforme)~: elle est à égalité pour le meilleur score d'apprentissage.
\end{solution}

\begin{solution}{exo:in:choosing:7}
Rétablies sur trois offres, les acceptabilités de premier rang deviennent 81.4~\% (teneur de marché), 13.5~\% (quant de banque) et 5.1~\%
(analyste de plateforme)~: retirer une offre dominée modifie les échelles et donc, un peu, les parts.
\end{solution}

\begin{solution}{exo:in:choosing:8}
Les pondérations sont tirées de façon uniforme par hypothèse, et non à partir d'une enquête auprès de personnes~; l'indice mesure la part de l'espace des pondérations qui favorise une offre,
ce qui ne dit rien du nombre de personnes qui retiennent telle ou telle pondération. Par ailleurs, les valeurs des attributs sont
illustratives.
\end{solution}

\begin{solution}{pb:in:choosing:1}
\begin{enumerate}
\item La rémunération et sa fourchette (13, 14), le lieu (15, 27), les horaires (26), le travail et ce qu'il mène à (16--25, 28), la sécurité (22).
\item Comme dans la définition du chapitre.
\item 1\,186\,354~\$, 1\,365\,186~\$, 903\,603~\$ et 845\,964~\$.
\item Afin que le poids accordé au risque soit explicite et puisse varier.
\item L'apprentissage, et les pondérations elles-mêmes.
\item Comme dans le tableau du chapitre.
\item Comme dans la définition~; l'offre de l'ingénieur en technologie.
\item 30, 15, 10, 15, 20, 10~; valeurs 0.852 (teneur de marché), 0.589 (plateforme), 0.584 (banque), 0.202 (technologie).
\item Comme dans la définition~; mise à l'échelle, criblage, pondération, poids d'échantillon sur le simplexe et dénombrement des rangs.
\item 79.5~\%, 15.9~\%, 4.6~\% et 0~\%.
\item De 0.452, en passant de 0.15 à 0.602.
\item Ce n'est le choix que pour quelqu'un qui accorde au risque de rémunération un poids supérieur à celui de tous les autres attributs réunis.
\item Des éléments sur ce qu'enseigne chaque poste~: sa \omterm{def:in:quant-researcher:feedback}{vitesse du retour d'information}, sa formation et ce que deviennent ses collaborateurs ensuite.
\item Il ferait monter à la fois la rémunération espérée et la fourchette de la plateforme.
\item Le tempérament, les personnes et la situation de la personne qui choisit.
\item Teneur de marché 79.5~\%, banque 15.9~\%, plateforme 4.6~\%, technologie 0~\%~; le poids du risque de rémunération doit augmenter de 0.452.
\item Dans un journal de décisions~: les attributs, les pondérations, les attentes et ce qui la ferait changer d'avis.
\item Lorsqu'une conviction change (rémunération, apprentissage, sécurité), et à une date fixée à l'avance, et non après un mauvais mois.
\item Un résultat qu'elle ne pouvait pas prévoir relève de la malchance~; un résultat qu'elle pouvait prévoir et ignorer relève d'un mauvais choix~; le
journal permet de les distinguer.
\item Sur les chiffres du chapitre, l'offre du teneur de marché l'emporte sous la plupart des pondérations, et seule une très forte aversion
au risque de rémunération favorise celle de la banque. Consigner pourquoi, et réexaminer lorsqu'une conviction, et non un résultat, change.
\end{enumerate}
\end{solution}

\begin{solution}{iq:in:choosing:1}
Nommez les attributs qui tranchent pour vous (le travail, ce que vous allez apprendre, les personnes) et montrez que vous connaissez les réalités du poste~;
ne pas dénigrer les autres offres.
\iqlookfor{raisons précises liées au poste.}
\end{solution}

\begin{solution}{iq:in:choosing:2}
Par le risque et le calendrier~: équivalents certains sous une aversion au risque déclarée, report et \omterm{def:in:how-pay-works:rsu}{déchéance} (ce que coûte un départ), et
calendrier fiscal.
\iqlookfor{la structure compte même à valeur espérée égale.}
\end{solution}

\begin{solution}{iq:in:choosing:3}
Comparer chaque couple~: l'offre $i$ est dominée s'il existe un $j$ au moins aussi bon partout et meilleur quelque part~; $O(n^2m)$ pour
$n$ offres et $m$ attributs, acceptable pour un petit $n$~; un algorithme de skyline fondé sur un tri fait mieux pour un grand $n$.
\iqlookfor{la définition par couples et son coût.}
\end{solution}

\begin{solution}{iq:in:choosing:4}
Tirer des variables exponentielles indépendantes et les normaliser (une loi de Dirichlet de paramètres unitaires), ou trier des tirages uniformes et
prendre les écarts~; normaliser des tirages uniformes ne donne pas une loi uniforme sur le simplexe.
\iqlookfor{la construction de Dirichlet et l'erreur courante.}
\end{solution}

\begin{solution}{iq:in:choosing:5}
Comparer le résultat à ce qui était connu et attendu au moment du choix (le journal)~: si le mauvais résultat tenait à un risque prévisible
que vous avez accepté en connaissance de cause, c'était de la malchance~; si vous avez ignoré une information dont vous disposiez, c'était une mauvaise décision.
\iqlookfor{éviter le biais du résultat.}
\end{solution}

\begin{solution}{iq:in:choosing:6}
Le modèle additif compte deux fois ce que les deux attributs ont en commun, ce qui leur donne un poids excessif~; les fusionner, réduire les poids
des deux, ou ne gardez que celui qui compte pour vous.
\iqlookfor{double comptage dans les modèles additifs.}
\end{solution}
